\documentclass[10pt,a4paper]{amsart}
\usepackage[margin=19mm]{geometry}
\usepackage{amsmath,amssymb,mathtools}
\usepackage[scr=rsfs,cal=euler]{mathalfa}
\usepackage{xfrac}
\usepackage[unicode,hidelinks,bookmarks=false]{hyperref}
\newcommand{\ie}{i.e.\ }
\newcommand{\cf}{cf.\ }
\newcommand{\resp}{respectively\ }
\newcommand{\Subsection}{\subsection}
\providecommand{\nobreakdash}{\nobreak\hbox{-}}
\setlength{\emergencystretch}{2em}
\multlinegap0pt
\usepackage{amssymb}
\usepackage{mathtools}
\usepackage{xcolor}
\usepackage[shortlabels]{enumitem}
\newtheorem{proposition}{Proposition}
\title{On the $C^2$-local systolic optimality of Zoll odd-symplectic forms}
\author{Samanyu Sanjay}
\date{Source: arXiv:2512.01937v3 (CC BY 4.0)}
\begin{document}
\maketitle

\noindent\emph{Source and licence.} On the $C^2$-local systolic optimality of Zoll odd-symplectic forms. Version arXiv:2512.01937v3 (CC BY 4.0), distributed by arXiv under the Creative Commons Attribution 4.0 International licence, http://creativecommons.org/licenses/by/4.0/, as recorded in the arXiv licence metadata for this exact version and shown on the abstract page https://arxiv.org/abs/2512.01937v3. Full licence notice: \texttt{licenses/arxiv-2512.01937-CC-BY-4.0.txt}.\par\smallskip

\noindent\textit{Editorial scope.} Counted unit: the article's proposition proposition (source0.tex lines 1004--1009). The article's complete original proof of it is delivered with it (source0.tex lines 1010--1090, one \texttt{proof} environment of 81 lines). No mathematics is written, repaired or re-derived: the statement and the proof are the article's own lines, in the article's own notation, and no retained line is substituted or re-flowed.  Retained uncounted as the source-local context the proof needs: lines 860--862; lines 971--973.  Nothing was omitted from any retained span, so the package carries no omission record.  No retained line contains a citation, so the wrapper carries no bibliography and no bibliography entry.  No retained line contains a figure environment, an image-inclusion command or drawing code, so no image is delivered and the package declares no asset dependency.  Editorial formatting only, in addition: an AMS \texttt{amsart} preamble that restates the article's own declarations and macros used by the retained lines, namely the environments \texttt{proposition} on separate counters, together with the standard packages those lines need.  The retained environments print in this document as Proposition 1 in order of appearance, whereas the article prints the same items with the numbering of its own full document.  The complete native source of the article as distributed in the arXiv e-print for this exact version is bundled unchanged as \texttt{sources/190169/source0.tex}.
\begin{align}\label{def_of_p_tilde}
    \tilde{p}^k_j \alpha_0 \wedge (\Omega_0)^{n-1} = \binom{n-1}{k} \binom{n-1-k}{j} d\big(\alpha_0 \wedge d\alpha^j_0 \wedge (\Omega_0)^{n-1-k-j}\wedge \eta \wedge d\eta^{k-1}\big)
\end{align}

\begin{align}\label{def_of_D}
    D:= \Tilde{P}+Q_1+Q_2
\end{align}

\begin{proposition}
Let $D: \Sigma \times\mathbb{R} \rightarrow \mathbb{R}$ be the function defined in \eqref{def_of_D} then, given an $\epsilon>0$ there exists a $\delta_0 >0$ such that the following is true when $s \in \mathbb{R}$ is close to $0$:
\begin{align}
 \vert\vert \Omega_0 - \Omega\vert\vert_{C^2} < \delta_0 \implies   \Big\vert\Big\vert\frac{\partial}{\partial s}D\Big\vert\Big\vert_{C^0} < \epsilon
\end{align}
\end{proposition}

\begin{proof}
\textbf{\underline{Step 1: The bounds on $p^k_j$}}\\
We need to show that given an $\epsilon>0$, there exist real constants $\delta, c>0$ such that $\vert\vert p^k_j\vert\vert_{C^0} < \epsilon$ if the following hold:
\begin{align}\label{est_bounds_basis_1}
  \max \bigg\{\vert\vert\eta\vert\vert_{C^0}, \vert\vert d\eta\vert\vert_{C^0}, \vert\vert\mathcal{F}\vert\vert_{C^0} \bigg\} < \delta \\\nonumber
\max \bigg\{\vert\vert\eta\vert\vert_{C^1}, \vert\vert d\eta\vert\vert_{C^1}, \vert\vert\mathcal{F}\vert\vert_{C^1} \bigg\} < c
  \end{align}
Since $\hat{\mathcal{F}}$ is the dual of $\mathcal{F}$, the bounds in \eqref{est_bounds_basis_1} imply that there exists a real constant $b_0 >0$ such that:
\begin{align}\label{F_dual_bounds}
    \vert\vert\hat{\mathcal{F}}\vert\vert_{C^0} < b_0 \delta \; \; \& \;\; \vert\vert\hat{\mathcal{F}}\vert\vert_{C^1} <  b_0 c
\end{align}
Now, recall that for each $1 \leq k \leq n-2$ and $0\leq j \leq n-1-k$ the functions $p^k_j$ satisfy:
\begin{align*}
     p^k_j \cdot \alpha_0 \wedge (\Omega_0)^{n-1} = \binom{n-2}{k} \binom{n-1-k}{j}  \Bigg[ d\hat{\mathcal{F}}(\alpha_0 \wedge d\alpha^{j+1}_0 \wedge \Omega^{n-1-k-j}_0 \wedge \eta \wedge d\eta^{k-2})\Bigg]
\end{align*}
Consider now the following differential form:
\begin{align}\label{def_of_gamma_k_j}
    \gamma_{k,j}:= \alpha_0 \wedge \eta \wedge d\alpha^j_0 \wedge \Omega^{n-1-k-j}_0 \wedge d\eta^{k-2}
\end{align}
The Libnitz formula along with \eqref{est_bounds_basis_1} implies that there exists a real number $b_1>0$ such that:
\begin{align}\label{gamma_c_0_bound}
    \vert\vert\gamma_{k,j}\vert\vert_{C^0} \leq b_1 \delta^{k-1} \; \;  \& \;\; \vert\vert\gamma_{k,j}\vert\vert_{C^1} \leq b_1 c^{k-1}
\end{align} 
Using the Leibnitz formula, \eqref{F_dual_bounds} and \eqref{gamma_c_0_bound} we see that the following is true:
\begin{align}\label{p_kj_bound_1}
    &\vert\vert d\hat{\mathcal{F}}[\alpha_0 \wedge \eta \wedge d\alpha^j_0 \wedge \Omega^{n-1-k-j}_0 \wedge d\eta^{k-2}]\vert\vert_{C^0} \\\nonumber \leq
     \vert\vert\hat{\mathcal{F}}&\vert\vert_{C^1} \cdot \vert\vert\gamma_{k,j}\vert\vert_{C^0} + \vert\vert\hat{\mathcal{F}}\vert\vert_{C^0} \vert\vert\gamma_{k,j}\vert\vert_{C^1} \leq b_0b_1 \delta^k + b_0b_1 c^k
\end{align} 
It is clear that by choosing $\delta, c$ in \eqref{est_bounds_basis_1} small enough we can make $\vert\vert p^k_{j}\vert\vert_{C^0} < \epsilon$ for any given $\epsilon >0$.\\
\\
\textbf{\underline{Step 2: The bounds on $\Tilde{p}^{k}_{j}$}}
\\
Recall from \eqref{def_of_p_tilde} that the functions $\Tilde{p}^k_j$ are defined by:
\begin{align*}
     \tilde{p}^k_j \alpha_0 \wedge (\Omega_0)^{n-1} = d\big(\alpha_0 \wedge d\alpha^j_0 \wedge (\Omega_0)^{n-1-k-j}\wedge \eta \wedge d\eta^{k-1}\big)
\end{align*}
In analogy with the previous step we make the following definition:
\begin{align}
    \hat{\gamma}_{k,j}:= \alpha_0 \wedge d\alpha^j_0 \wedge (\Omega_0)^{n-1-k-j}\wedge \eta \wedge d\eta^{k-1}
\end{align}
Now, it is clear from \eqref{est_bounds_basis_1} that there exists a real $d_1 >0$ such that:
\begin{align}
    \vert\vert\hat{\gamma}_{k,j}\vert\vert_{C^1} < d_1c^k
\end{align}
So, we now conclude that the following holds for each $1 \leq k \leq n-2$ and $0\leq j \leq n-1-k$:
\begin{align}
\vert\vert \tilde{p}^k_j\vert\vert_{C^0} \leq  \vert\vert\hat{\gamma}_{k,j}\vert\vert_{C^1} <  d_1c^{k}
\end{align}
\\
\textbf{\underline{Step 3: The bounds on $\tilde{Q}^{k,j}_1$}}
\\
First recall that the functions $\tilde{Q}^{k,j}_1$ are defined by the equations:
\begin{align}
         \tilde{Q}^{k,j}_1 \cdot \alpha_0 \wedge (\Omega_0)^{n-1} = d\hat{\mathcal{F}}[\alpha_0 \wedge \eta \wedge (\Omega_0)^{n-1-k-j} \wedge (d\alpha_0)^{j} \wedge d\eta^{k-1}]
\end{align}
So, if we write $\kappa_{k,j}:= \alpha_0 \wedge \eta \wedge (\Omega_0)^{n-1-k-j} \wedge (d\alpha_0)^{j} \wedge d\eta^{k-1}$. Then, by \eqref{est_bounds_basis_1} there exists a real constant $d_2>0$ such that the following hold:
\begin{align}
    \vert\vert\kappa_{k,j}\vert\vert_{C^0} \leq d_2 \delta^{k} \; \;  \& \;\; \vert\vert\kappa_{k,j}\vert\vert_{C^1} \leq d_2 c^{k}
\end{align}
This allows us to conclude as in step 1 above that there exists a $d_3>0$ such that:
\begin{align}
    \vert\vert\tilde{Q}^{k,j}_1\vert\vert_{C^0} \leq \vert\vert\hat{\mathcal{F}}\vert\vert_{C^0}\cdot \vert\vert\kappa_{k,j}\vert\vert_{C^1} + \vert\vert\hat{\mathcal{F}}\vert\vert_{C^1}\cdot \vert\vert\kappa_{k,j}\vert\vert_{C^0} \leq d_3(\delta^{k+1} + c^{k+1})
\end{align}
\\
\textbf{\underline{Step 4: The bounds on $\Tilde{Q}^{k,j}_2$}}
\\
\\
First, recall that for each $k$ and $j$ the functions $\Tilde{Q}^{k,j}_2$ are defined by the equations:
\begin{align}
        \tilde{Q}^{k,j}_2 \cdot \alpha_0 \wedge (\Omega_0)^{n-1} = d( \alpha_0 \wedge \eta  \wedge (\Omega_0)^{n-1-k-j} \wedge (d\alpha_0)^j \wedge  (d\eta)^{k})
    \end{align}

\begin{align}
    \vert\vert\tilde{Q}^{k,j}_2\vert\vert_{C^0} \leq a_1c^{k+1}
\end{align}
So, by choosing $c$ in \eqref{est_bounds_basis_1} small enough we can guarantee that for a given $\epsilon>0$ and for each $k$ and $j$ the functions $\Tilde{Q}^{k,j}_2$ satisfy the bounds:
\begin{align}
    \vert\vert\tilde{Q}^{k,j}_2\vert\vert_{C^0} < \epsilon
\end{align}
Since $D(x,s)$ is a polynomial in $s$ with coefficients given by $\tilde{P}$, $Q_1$ and $Q_2$ we see that if $\delta$ and $c$ in \eqref{est_bounds_basis_1} are small enough then $\vert\vert\frac{\partial D}{\partial s}\vert\vert_{C^0} < \epsilon$ as claimed. 
\end{proof}

\end{document}
